%% ma: L10-B14-0052
%% lop: 10
%% chuong: 5
%% bai: 14
%% dang: 10.14.9
%% loai: TL
%% muc: TH
%% dap_an: "-"
%% nguon: "chuyen-de-cac-so-dac-trung-cua-mau-so-lieu-khong-ghep-nhom-toan-10-chuong-trinh-moi.pdf (GV. Phan Nhật Linh) (P05), Bài 03 Dạng 2 Bài tập tự luận 12"
%% kien_thuc: "Lớp 10 Bài 13 (số trung bình), Bài 14 (phương sai, độ lệch chuẩn)"
%% trang_thai: chinh-thuc
%% ghi_chu: "Tự duyệt theo yêu cầu giáo viên 22/10/2026."
\begin{ex}
Kết quả $5$ bài kiểm tra môn Toán (thang điểm 10) của hai bạn Lan và Hoa như sau:
\begin{center}
\begin{tabular}{|c|c|c|c|c|c|}\hline
Lan & $4$ & $7$ & $6$ & $8$ & $9$\\\hline
Hoa & $7$ & $6$ & $8$ & $7$ & $6$\\\hline
\end{tabular}
\end{center}
Bạn nào có kết quả kiểm tra môn Toán đồng đều hơn?
\loigiai{$\overline x_L=\overline x_H=6{,}8$.

$s_L^2=\dfrac{(4-6{,}8)^2+(7-6{,}8)^2+(6-6{,}8)^2+(8-6{,}8)^2+(9-6{,}8)^2}5=2{,}96$; $s_H^2=0{,}56$.

Vì $s_H^2<s_L^2$ nên Hoa có kết quả đồng đều hơn.}
\end{ex}
